\section{Matrix}
TODO

\subsection{Matrix Properties}
TODO

\subsection{Matrix Operations}
TODO

\subsection{Matrix as a Combination of Vectors}
TODO

\subsection{Matrix as a Representation of Systems of Linear Equations}
TODO

\subsection{Matrix as a Means of Linear Transformation}
TODO: explain Rotation, Scale, Shearing...

\subsection{Why is Translation not a form of Linear Transformation}
TODO: explain why translation is actually Shearing from a higher dimension projected down
to a dimension lower as translation...

\subsection{Matrix for Transforming Between Two Eucledian Spaces}
TODO

\subsection{Projection Matricies}
TODO

\subsection{Useful Matrix Constructs}
You can rotate an arbitrary point $A$ around another point $O$ by following these steps:
\begin{enumerate}
    \item Translate $O$ back to the origin, and bring $A$ with it.
    \item Rotate this new $A$ around the origin(which is also where the new $O$ resides)
        as how you would normally.
    \item Translate the new $O$ back to where $O$ originally was, and bring the new rotated $A$ with it.
    \item You now have your rotated $A$ around $O$!
\end{enumerate}

This can be further written in matrix if we promote $O$ and $A$ to homogeneous coordinates(to allow for translation),
in which the matrix $\textbf{T}$ represents translating $O$ back to origin, and matrix $\textbf{R}$ represents
some pure rotational matrix around the origin. Which can be written mathematically as:
\begin{center}
    $\textbf{M}=\textbf{T}^{-1}\textbf{R}\textbf{T}$
\end{center}

In which you can now apply the final matrix $\textbf{M}$ on any point $A$ which effectively rotates it around $O$!

\sgap
You can create a matrix to convert from one Eucledian space $A$ to Eucledian space $B$ by filling a matrix with
basis vectors per column of space $A$ relative to $B$, in mathematical terms:

TODO



